\chapter{Introduction}
\label{ch:introduction}

% Budget and sources: docs/thesis/outline.md. Every claim about behaviour points at a test, a figure or a number in data/.

\section{The problem}
\label{sec:intro-problem}

A payment dispute begins with a customer who looks at a statement and objects to a line on it. The
charge may be one they do not recognise at all, a purchase that never arrived, a subscription
taken twice, or an amount larger than the one they agreed to. Whatever the reason, the customer
does not take it to the shop. They take it to their own bank, because that is where the money
left from, and because the law in most jurisdictions gives them rights against that bank that it
does not give them against the shop.

For the bank the complaint is the start of a process with several strands running at once. There
are legal deadlines, which differ by jurisdiction and payment method and some of which are
measured in business days. There is a duty, in several regimes, to give the customer their money
back before anyone knows whether the complaint is justified. There is a second process on the
other side of the card network, in which the bank tries to recover that money from the merchant's
bank and the merchant answers with evidence. There are letters and emails the customer is owed at
particular points, with particular content. And there is the money itself, which must move in the
bank's own core system exactly once, and move back if the law allows and the outcome requires it.
Each of these strands is simple in isolation. The difficulty is that they all depend on the same
question, which is where the dispute currently stands, and that a mistake in any one of them is
either a breach of a legal duty or a loss.

The strands are complicated further by the customer who disputes a payment they did in fact make.
First-party fraud, often called friendly fraud, is a dispute raised in bad faith, or at least in
error, by the legitimate cardholder: the parcel that did arrive, the purchase a family member made
with the card, the subscription the customer forgot they had started. The bank cannot treat every
customer as a suspect; the regulation that protects customers assumes good faith and frequently
requires the refund to come first. Nor can it ignore that some disputes are not what they claim to
be. The position this thesis takes is in its title. The customer is given the benefit of the doubt
at the start, with the refund or credit the law requires, and the evidence then decides how the
dispute ends. The engine's job is to make both halves of that sentence hold at once.

\section{Why disputes are regime-specific}
\label{sec:intro-regimes}

The rules a bank follows depend on the payment rail and the jurisdiction, and the differences
between them are not merely different numbers in the same process. Under the European Payment
Services Directive, an unauthorised card payment is refunded by the end of the next business day
and the refund is final unless the bank later recovers it \cite{psd2}. A SEPA direct debit can be
claimed back within eight weeks with no reason given and no defence available to the merchant
\cite{sepacore}. The United States' Regulation~E, for debit cards, instead requires a provisional
credit while the bank investigates \cite{rege}, and Regulation~Z, for credit cards, issues no
credit at all and suspends the collection of the disputed amount \cite{regz}. Some regimes have a
merchant on the other side who may contest the claim; one has nobody. Chapter~\ref{ch:background}
sets these regimes out in detail.

Software that handles disputes under more than one regime therefore has a choice to make. It can
treat each regime as its own workflow, with its own states and its own code, and accept that a
fix in one is not a fix in the others. Or it can look for the lifecycle the regimes share and
express their differences as configuration of that one lifecycle. This thesis argues for the
second: a single finite state machine, parameterised by the regulatory regime, in which the
regime decides which states are reachable, which deadlines run and which money moves, and the
state machine itself is the same for all of them. The system is written from the issuer's side,
that is, from the position of the customer's own bank, which receives the complaint, owes the
legal duties and files the chargeback outward. Two European regimes, PSD2 card payments and SEPA
direct debits, are the focus of the work; two American regimes are modelled alongside them to
show that the model generalises rather than to serve a market.

\section{What was built}
\label{sec:intro-built}

The result is a dispute engine that a bank or payment provider can use as a service, with an
analyst workbench in front of it. Several institutions share one deployment, and the database
itself keeps each one's data apart. The engine models the dispute lifecycle as a state machine
gated by regime, starts the regulatory clocks each regime requires in the institution's own
business-day calendar, posts every movement of money to a ledger and through a port to the
institution's core banking system, collects the customer's account of what happened through a
questionnaire chosen by the dispute's reason, sends the letters and emails the customer is owed,
and scores each dispute with readable rules that can hold a credit until a person justifies it.
Who may do what inside an institution is decided by policy over the authenticated caller, and
the caller is recorded on every event. Chapter~\ref{ch:design} takes these parts in turn.

The work aims to be more than a prototype in one respect in particular: its claims can be
checked. The behaviour of each regime is pinned by tests that walk every transition the regime
allows and refuse the ones it forbids; isolation between institutions is tested against the
database rather than assumed from the code; and every number in the evaluation comes from a
measurement run that is committed to the repository and can be repeated with one command
(chapter~\ref{ch:evaluation}).

\section{Structure of the thesis}
\label{sec:intro-structure}

Chapter~\ref{ch:background} introduces the parties to a dispute and the four regulatory regimes,
and draws from them what a dispute engine must guarantee. Chapter~\ref{ch:requirements} turns
those guarantees into requirements, lists the use cases, and records what was deliberately left
out. Chapter~\ref{ch:design} presents the design, one decision at a time: the regime-parameterised
state machine, the clocks, the ledger and the banking core, the questionnaire, communications and
fraud scoring, and the tenancy and authorization model. Chapter~\ref{ch:implementation} describes
how the design was built and operated, from the generated API contract to the analyst workbench
and its observability. Chapter~\ref{ch:evaluation} evaluates correctness, isolation, concurrency
and performance, and states the limitations. Chapter~\ref{ch:conclusion} concludes and outlines
future work. The appendix contains excerpts of the API contract and the raw evaluation tables.
